functions to document:

X playgesture

X playGestureFromSequencerStep()

X midiInput

X updateStepRow()

x monophonic-sequencer 1697-1782 (not documented in thesis.tex)

x polyphonic-sequencer 1787-1848 (not documented in thesis.tex)

x calculateEuclideanDistance() 1915-1959 - not doc'd

calculateCentrality()  1960 - 1999

\chapter{Patch History Window Code Excerpts}
\label{app:patchHistoryExcerpts}

\begin{table}[h]
  \centering
  \caption{Overview of Notable Code Listings from patchHistory.js, the script for the Patch History Window}
  \begin{tabular}{@{}llll@{}}
    \toprule
    \textbf{Listing} & \textbf{Function} & \textbf{Purpose} & \textbf{Lines} \\
    \midrule
        \ref{lst:playGesture-function} & playGesture() & Gesture playback from selected changeNode & 2370-2459 \\
        \ref{lst:playGestureFromSequencerStep-function} & playGestureFromSequencerStep() & Gesture playback from sequencer & 2022-2067 \\
        \ref{lst:midi-input} & MIDI CC Input & Map MIDI CC to graph traversal & 283-315 \\
        \ref{lst:updateStepRow} & updateStepRow() & Assign a changeNode to a step & 583-660 \\
        \ref{lst:monophonic-sequencer} & Monophonic Playback & Steps recalled one at a time & 1697-1782 \\
        \ref{lst:polyphonic-sequencer} & Polyphonic Playback & Each step becomes its own loop & 1787-1848 \\
        \ref{lst:calculateEuclideanDistance} & calculateEuclideanDistance() & Map changeNode distance onto step length & 1915-1959 \\
        \ref{lst:calculateCentrality} & calculateCentrality() & Map changeNode centrality onto step length & 1960 - 1999 \\

    \bottomrule
  \end{tabular}
\end{table}



\section{playGesture()}

\begin{lstlisting}[language=JavaScript, caption={Plays a gesture when the player either selects its associated changeNode, or from within the Gesture Editor}, label={lst:playgesture-function}]]

function playGesture(mode) {
        gestureCy.nodes().removeClass('highlighted')
        // 'repeat' is passed by the function call when looping is on, so we don't want to have to get the same data again if the loop is on
        if(mode != 'repeat'){
            // reset the gesture scheduler
            gestureData.scheduler = [ ]
            
            // sort objects by timestamp
            // sortedGestureNodes = [...gestureData.nodes].sort((a, b) => a.data.timestamp - b.data.timestamp);
            // Get the starting timestamp (the earliest one)
            // gestureData.startTime = sortedGestureNodes[0].data.timestamp;
            // gestureData.endTime = sortedGestureNodes[sortedGestureNodes.length - 1].data.timestamp;
            // gestureData.length = gestureData.endTime - gestureData.startTime
        }
 
        // create the scheduler
        gestureData.nodes.forEach((node) => {
            const delay = node.data.timestamp - gestureData.startTime; // Calculate delay from the start

            // Use setTimeout to schedule the callback
            const timeoutID = setTimeout(() => {
                
                // highlight the node
                let hNode = gestureCy.getElementById(node.data.id);
                hNode.addClass('highlighted');

                if(gestureData.assign.param === 'default'){
                    
                    let data = {
                        parent: node.data.parents,
                        param: node.data.param,
                        value: node.data.value
                    }
                    sendToMainApp({
                        cmd: 'playGesture',
                        data: data,
                        kind: 'n/a'
                    })

            
                } else {
                    // process it using the gesturedata assign range data for scaling

                    // convert the value from the source value's min and max to gestureData.assign.range
                    // first get the min and max of the source value
                    // synthParamRanges        
                    let value = node.data.value
                    
                    let storedParam = patchHistory.synthFile.audioGraph.modules[node.data.parents].moduleSpec.parameters[node.data.param]
                    let targetParam = gestureData.assign
                    
                    sendToMainApp({
                        cmd: 'playGesture',
                        data: convertParams(storedParam, targetParam, value),
                        kind: targetParam.kind

                    })
            
                }

                if(UI.gestureEditor.loop.checked && gestureData.length === delay){
                    playGesture('repeat')
                    // setTimeout(() => {
                    //     playGesture('repeat')
                    // }, 250);
                }
            }, delay);



            gestureData.scheduler.push(timeoutID)
        });

        // get the end of the gesture
        const maxDelay = Math.max(...gestureData.nodes.map(node => node.data.timestamp - gestureData.startTime));

        const finalTimeoutID = setTimeout(() => {
            gestureCy.elements().removeClass('highlighted');

            // if looping is off, set the stop button back to 'play'
            if (!UI.gestureEditor.loop.checked) {
                UI.gestureEditor.play.textContent = 'Play'
            }
        }, maxDelay + 1); // +1 to ensure it's last

        gestureData.scheduler.push(finalTimeoutID);
        


    }
\end{lstlisting}

\section{functionName}

\begin{lstlisting}[language=JavaScript, caption={This function plays a gesture when triggered from a sequencer step. It uses the quantizeGesture function to ensure that it plays back entirely within the current step length.}, label={lst:playGestureFromSequencerStep-function}]]

// playback a stored gesture from a sequencer step
    function playGestureFromSequencerStep(gesture, stepLength){
        
        let quantizedGesture = quantizeGesture(gesture, stepLength)
        
        // create the scheduler
        quantizedGesture.forEach((node) => {
            const delay = node.t * 1000; // (convert to milliseconds)
            
            // Use setTimeout to schedule the callback
            const timeoutID = setTimeout(() => {
                if(gesture.assign.param === 'default'){
                    
                    let data = {
                        parent: node.parent,
                        param: node.param,
                        value: node.value
                    }
                
                    sendToMainApp({
                        cmd: 'playGesture',
                        data: data,
                        kind: 'n/a'
                    })
    
                    
                } else {

                    let value = node.value
                    
                    let storedParam = patchHistory.synthFile.audioGraph.modules[node.parent].moduleSpec.parameters[node.param]
                    let targetParam = gesture.assign
                    
                    sendToMainApp({
                        cmd: 'playGesture',
                        data: convertParams(storedParam, targetParam, value),
                        kind: targetParam.kind
    
                    })
            
                }
    
            }, delay);
    
            gesture.scheduler.push(timeoutID)
        });
    }
    
    function quantizeGesture(gesture, stepLength) {
        const duration = gesture.endTime - gesture.startTime;
        const scale = stepLength / duration;
      
        // Map each point's timestamp to the new interval [0, stepLength]
        return gesture.gesturePoints.map(point => ({
          ...point,
          t: (point.timestamp - gesture.startTime) * scale
        }));
    }

\end{lstlisting}

\section{functionName}

\begin{lstlisting}[language=JavaScript, caption={The createMerge function takes two changeNodes selected by the player and merges them. A new branch is created to store the result of the merge. The unified patch state is then loaded into the main application and rendered in the visual and audio synth graphs.}, label={lst:functionName-function}]]

% INSERT A FUNCTION FROM PATCHHISTORY.JS HERE 
\end{lstlisting}

\bigskip
\textbf{Complete File} — The full, commented source is archived at
\\https://github.com/michaelpalumbo/forkingpaths/blob/dissertationJune/src/scripts/patchHistory.js
